\section{Checkable hypotheses and their limits}
\label{sec:hypotheses}

\subsection{Stable convexity of the quadratic parts}

A tuple $(A_1,\ldots,A_m)$ is \emph{stably convex} if it has a
neighborhood in $(\Sym^n)^m$ such that every tuple $(M_1,\ldots,M_m)$
in that neighborhood has a convex joint image
$\{(x\transpose M_i x)_i:x\in\R^n\}$. This is stable convexity in
the sense of \citet{Sheriff2013}, rather than convexity of each
component function.

\begin{proposition}\label{prop:stable}
If the quadratic-part tuple $(A_i)_i$ is stably convex, then $f^h$
has asymptotic hyperplane convexity, for arbitrary $b_i,c_i$.
Consequently the certificate equivalence holds whenever $S\ne\varnothing$
and either
\begin{enumerate}
\item $m=2$, in any dimension $n\geq1$; or
\item $m=3$, $n\geq3$, and some real combination
      $\sum_i\theta_i A_i$ is positive definite.
\end{enumerate}
\end{proposition}
\begin{proof}
For $s\ne0$, parameterizing $H_{\alpha,s}$ by
$(x,\alpha\transpose x/s)$ makes its image the joint image of
\begin{equation}\label{eq:perturbed-parts}
 M_i(s)=A_i+s^{-1}(\alpha b_i\transpose+b_i\alpha\transpose)
                         +s^{-2}c_i\alpha\alpha\transpose.
\end{equation}
For each fixed $\alpha$, these matrices tend to $A_i$, so stable
convexity gives convexity for all sufficiently large $s$.
Dines' theorem gives convexity for every two-form tuple, including
every perturbation and hyperplane restriction.
For three forms on $\R^n$, $n\geq3$, Polyak's theorem states that
a positive definite real combination implies convexity of the joint
image \citep[Theorem~2.1]{Polyak1998}. Positive definiteness persists
under small perturbations, so this gives stable convexity and the
claimed conclusion.
\end{proof}

The first case is the known two-constraint setting. The second is a
convenient criterion on the quadratic blocks, not a new three-form
convexity theorem. One must distinguish positive
definiteness in $\spanop\{A_i\}$ from positive definiteness in
$\spanop\{Q_i\}$. The latter implies the former, but in the delicate
case of a trivial certificate they are equivalent. Indeed, if
$\lambda^0\geq0$, $\lambda^0\ne0$, and
$A_{\lambda^0}=b_{\lambda^0}=0$, strict feasibility gives
$c_{\lambda^0}<0$. Thus
$e_{n+1}e_{n+1}\transpose=Q_{\lambda^0}/c_{\lambda^0}$ belongs to
the real span of the $Q_i$. If
$\sum_i\theta_iQ_i=\left(\begin{smallmatrix}A&b\\b\transpose&c\end{smallmatrix}\right)$
has $A\succ0$, adding $r e_{n+1}e_{n+1}\transpose$ with
$r>b\transpose A^{-1}b-c$ makes it positive definite by the Schur
complement. If no nonzero trivial certificate exists, either a
nontrivial one already exists or $\mathcal K=\{0\}$; under convexity
of the quadratic-part image, the latter case was covered in
\citet[Proposition~2.14]{BDSpreprint}. Thus the block formulation
resolves no previously open PDLC case beyond the main certificate
theorem.

\begin{example}[Asymptotic hyperplane convexity is strictly weaker than HHC]
\label{ex:ahc-not-hhc}
In $\R^3$, set
\[
 \begin{split}
 f_1&=x_1^2+x_2^2+x_3^2-1,\\
 f_2&=x_1^2-x_2^2+x_3^2-1,\\
 f_3&=2x_1x_2+x_3^2-1.
 \end{split}
\]
The origin is strictly feasible. Since $A_1=I$, the tuple is stably
convex by the three-form result, and Proposition~\ref{prop:stable}
gives asymptotic hyperplane convexity.
On the homogeneous hyperplane $x_3=t$, however, the image is
\[
 \{(x_1^2+x_2^2,\ x_1^2-x_2^2,\ 2x_1x_2):x_1,x_2\in\R\}.
\]
Every image point $(y_1,y_2,y_3)$ satisfies
$y_1^2=y_2^2+y_3^2$. The image contains $(1,1,0)$ and $(1,-1,0)$
but not their midpoint $(1,0,0)$. Thus HHC fails. This also proves
that stable convexity of the quadratic parts does not imply HHC of
the homogenization.
\end{example}

\begin{example}[HHC need not give stable convexity of the parts]
\label{ex:hhc-not-stable}
Let $f_1=f_2=f_3=x_1^2-x_2^2-1$ in $\R^3$. On every homogeneous
hyperplane the three components agree. A single real homogeneous
quadratic form has image $\{0\}$, a ray, or $\R$, as follows by
scaling and, in the indefinite case, evaluating both signs. Hence
every such hyperplane image is convex.
For any $\varepsilon>0$, perturb the three quadratic parts to
\[
 u=x_1^2-x_2^2,\qquad
 u+\varepsilon(x_1^2+x_2^2),\qquad
 u+2\varepsilon x_1x_2.
\]
Their values at $e_1,e_2$ have midpoint $(0,\varepsilon,0)$.
A preimage would require $x_1^2=x_2^2=1/2$ and $x_1x_2=0$, which
is impossible. These arbitrarily small perturbations are nonconvex.
Together with Example~\ref{ex:ahc-not-hhc}, this separates the two
sufficient conditions in both directions.
\end{example}

\subsection{Why ordinary hidden convexity is insufficient}

\begin{example}\label{ex:ordinary-hc}
Consider, in $\R^3$ with $x_3$ free,
\begin{equation}\label{eq:ordinary-hc-example}
 f_1=x_1^2-x_2^2-1,\quad f_2=x_2^2-x_1^2-1,\quad
 f_3=x_1x_2-1,\quad f_4=-x_1x_2-1.
\end{equation}
The strict set contains the origin. On $S$,
\[
 (x_1^2+x_2^2)^2=(x_1^2-x_2^2)^2+4x_1^2x_2^2<5.
\]
It lies in the convex cylinder
$\{x:x_1^2+x_2^2<\sqrt5\}$, so its hull is proper.
Every aggregated quadratic matrix is traceless and every linear
part is zero. Positive semidefiniteness therefore forces the matrix
to vanish, so all convex certificates are trivial. More explicitly,
the only conditions are $\lambda_1=\lambda_2$ and
$\lambda_3=\lambda_4$.

Nevertheless both the full homogeneous image and the image on $E$
are convex. Put $u=x_1^2-x_2^2$, $v=x_1x_2$, and $w=t^2$.
The complex square map shows that $(u,v)$ ranges over $\R^2$;
$w$ ranges independently over $\R_+$. Hence the full image is
the linear image of $\R^2\times\R_+$ under
\[
 (u,v,w)\longmapsto(u-w,-u-w,v-w,-v-w).
\]
The image on $E$ is its restriction to $w=0$ and is convex too.

For $s\ne0$ restrict instead to $x_1=st$. The same linear map is
injective, and $w=x_1^2/s^2$. The two preimages
$(x_1,x_2,x_3,t)=(1,0,0,1/s)$ and $(0,1,0,0)$ give a midpoint
whose $(u,v,w)$ coordinates would have to be
$(0,0,1/(2s^2))$. They would force
$x_1^2=x_2^2=1/2$ and $x_1x_2=0$, a contradiction.
Thus no nonzero level in this family has a convex image. This is
the precise obstruction to asymptotic hyperplane convexity.
Proposition~\ref{prop:shor-closure} also gives $P_{\rm SDP}=\R^3$ here,
although the true hull is bounded in two coordinates.

Three inequalities already suffice for a proper-hull obstruction:
delete $f_4$. The trace argument and both ordinary hidden-convexity
statements remain valid. With $p=x_1+x_2$ and $q=x_1-x_2$, strict
feasibility requires $|pq|<1$ and $p^2-q^2<4$.
If $|p|\geq3$, then $|q|<1/3$ and
$p^2<4+1/9$, a contradiction. Thus $|p|<3$ is a proper affine
bound. The set is unbounded, since $x_1=-x_2$ is feasible for every
value of $x_1$. The three-coordinate map
$(u,v,w)\mapsto(u-w,-u-w,v-w)$ is also injective, so the preceding
hyperplane obstruction persists.
\end{example}

\subsection{The role of strict feasibility for closed systems}

\begin{example}\label{ex:closed}
In $\R^3$, let
\[
 f_1=x_1x_2-1,\qquad f_2=-x_1x_2+1,\qquad f_3=x_1^2-x_2^2.
\]
Then $S=\varnothing$, whereas
\[
 T=\{x:x_1x_2=1,\ |x_1|\leq|x_2|\}
\]
is nonempty and has $|x_1|\leq1$. Consequently its hull is proper.
The convex hull of the four feasible planar points
$\pm(1,1)$, $\pm(1/2,2)$ contains a neighborhood of the origin:
the two positive points are linearly independent and their symmetric
convex hull is a nondegenerate parallelogram. Since $x_3$ is free,
$0\in G:=\interior\conv T$ and $G$ has full dimension.

The quadratic matrices are traceless and the linear parts vanish,
so every convex certificate is trivial. HHC holds because
$Q_2=-Q_1$; every hyperplane image is a linear image of a two-form
image, which is convex by Dines' theorem.

This example also shows that the condition
$Q_\lambda\ne0$ for all nonzero $\lambda\geq0$ in
\citet[Theorem~2.24 and Remark~2.25]{BDSpreprint} cannot simply
be removed. In their notation a good closed-system multiplier
must satisfy $G\subseteq\{f_\lambda<0\}$ and give $Q_\lambda$
at most one negative eigenvalue. Since $0\in G$, any such multiplier
would have $c_\lambda=\lambda_2-\lambda_1<0$.
Put $\mu=\lambda_1-\lambda_2>0$. Its homogeneous matrix is
\[
 Q_\lambda=\operatorname{diag}\left(
  \begin{pmatrix}\lambda_3&\mu/2\\\mu/2&-\lambda_3\end{pmatrix},
   0,-\mu\right).
\]
The $2\times2$ block has determinant
$-\lambda_3^2-\mu^2/4<0$, so it contributes one negative
eigenvalue; $-\mu$ contributes a second. There are therefore no
good closed-system multipliers, whereas
$\varnothing\ne G\ne\R^3$. Also $Q_{(1,1,0)}=0$.
Thus even taking an interior of the intersection of the good
aggregations would give the whole space, not $G$.
\end{example}

\subsection{A certificate need not describe the hull}

\begin{example}\label{ex:strip}
In $\R^3$ with $x_2,x_3$ free, let
\[
 f_1=x_1^2-4,\qquad f_2=-x_1^2+2x_1+3.
\]
Dines' theorem gives HHC. Factoring the two inequalities gives
$S=\{x:-2<x_1<-1\}$, which is already convex.
Globally convex nonzero aggregations have
$\lambda_1\geq\lambda_2\geq0$, with $\lambda_1>0$, and hence
\[
 f_\lambda(0)=-4\lambda_1+3\lambda_2<0.
\]
Thus even all strict globally convex aggregations admit the origin,
although $x_1<-1$ is valid for $S$. The Shor projection also admits
$x=0$, with $X=\operatorname{diag}(7/2,0,0)$: the two lifted
residuals both equal $-1/2$.
The relaxation is still proper, since the first convex quadratic
inequality confines $|x_1|\leq2$. Therefore the whole-space
equivalence does not give hull exactness.
\end{example}

These examples separate the hypotheses, the open and closed systems,
and the strength of the resulting certificate. Proposition~\ref{prop:stable}
supplies concrete classes satisfying the hypotheses; the certificate
theorem gives no general procedure for verifying HHC or asymptotic
hyperplane convexity.
